%% NeurIPS 2026 Workshop Paper Checklist
%% Interpretability for Discovery Workshop

\section*{Paper Checklist}

\begin{enumerate}

\item \textbf{Claims.}
Do the main claims made in the abstract and introduction accurately reflect the paper's contributions and scope?
\answerYes{} The abstract and introduction present three specific claims (circuit structure, period-3 carrier, faithfulness) that are directly supported by the results in §3. Each claim is scoped to one checkpoint (SpliceBERT.1024nt), one task (masked donor prediction), and one behavioral assay (canonical short introns). No transfer to biological mechanisms or other models is asserted beyond what is measured.

\item \textbf{Limitations.}
Did you discuss the limitations of your work?
\answerYes{} §4 (Discussion) explicitly identifies three open questions: (a) the acceptor-payload semantic identity is unresolved and cannot be named on the current checkpoint; (b) the $\approx$10.63-nt carrier's OV function in the donor circuit is bounded at sub-SESOI levels and its general function is open; (c) cross-example specificity of the acceptor payload is not separable at current precision. §3.3 clarifies that the named circuit is a structured core within an irreducibly distributed computation, not a minimal sparse generator.

\item \textbf{Theory assumptions and proofs.}
If the paper includes theoretical contributions, do all theorems, corollaries, etc.\ appear to be correct?
\answerNA{} This paper contains no theoretical claims or proofs.

\item \textbf{Experimental reproducibility.}
Are the experimental results reproducible?
\answerYes{} All results are derived from the published SpliceBERT checkpoint (publicly available) and GENCODE v43 annotations (public). Code, preregistrations, and JSON result artifacts are included in the supplementary repository. Every claim is controlled by an immutable artifact file listed in the results manifest. Pre-registered protocols with smoke-gate tests and adversarial verification were used throughout.

\item \textbf{Code.}
Is code available, or will it be made available?
\answerYes{} The full analysis codebase will be released upon de-anonymization. The codebase includes all runner scripts, pre-registration documents, and result manifests.

\item \textbf{Compute.}
Does the paper include the compute required to reproduce the results?
\answerYes{} All experiments use a single pre-trained SpliceBERT checkpoint (inference only, no training). Compute requirements are modest: all analyses run on a single GPU in under 48 hours total. No large-scale compute is required to reproduce these results.

\item \textbf{Data.}
Does the paper reference all data used?
\answerYes{} The model (SpliceBERT.1024nt) is publicly available~\citep{chen2024splicebert}. Genomic annotations use GENCODE v43~\citep{mudge2023gencode}. RNA-FM and RNA-BERT checkpoints are publicly available~\citep{chen2022rnafm,akiyama2022rnabert}. The human-only SpliceBERT sibling checkpoint is publicly available from the same source as SpliceBERT.

\item \textbf{IRB approvals.}
If the paper involves human subjects, was ethical approval obtained?
\answerNA{} This paper involves no human subjects. All data are publicly available genomic annotations and pre-trained model weights.

\item \textbf{Responsible use.}
Does the paper discuss potential positive societal impacts, as well as any potential negative impacts?
\answerYes{} See below.

\end{enumerate}

\subsection*{Responsible Use Statement}

\paragraph{Potential positive societal impacts.}
This work develops mechanistic interpretability methods applicable to genomic foundation models. The methods — reconstruct--inject patching, position-table Fourier analysis, and nucleotide-invariant causal surgery — are designed to surface internal computations that are not accessible by behavioral probing alone. Positive impacts include: (1) improved understanding of what genomic models learn from multi-species pretraining data; (2) a transferable methodological template for auditing other GFMs for spurious or biologically meaningful features; (3) identification of open questions (e.g.\ acceptor-payload identity) that can guide future model evaluation and data collection.

\paragraph{Potential risks and mitigations.}
The methods described in this paper are interpretability tools applied to an existing published model; they do not enable training new genomic models or generating novel sequences. The primary risk is misinterpretation of results: (1) the donor--acceptor circuit is a model mechanism, not a direct claim about spliceosome biochemistry or transcript abundance; no external biological assay is used; (2) the period-3 carrier is attributed to multi-species pretraining by comparative analysis, not by retraining — origin dynamics cannot be recovered from a single checkpoint; (3) the faithfulness results show the named circuit is not a minimal sparse generator, and minimality is explicitly unclaimed. We have taken care throughout to scope claims precisely (see §3 and §4) and to avoid overstating biological implications. The open-source release of code and preregistration documents will allow independent verification and replication.

\paragraph{Broader context.}
Genomic foundation models are increasingly used in research and clinical genomics contexts. Mechanistic interpretability provides a tool for auditing these models before deployment, identifying features that may be spurious artifacts of training data rather than reflections of underlying biology. We view this work as a step toward model transparency in a domain where the stakes of misinterpreted predictions are high.
